\onecolumn

\section{Fase 2: Análisis y selección de literatura}

En esta fase se realizó una lectura estratégica de los artículos seleccionados, enfocándose en el \textit{abstract}, la introducción y las conclusiones. El objetivo fue identificar información clave que permita construir el estado del arte del problema de investigación.

Para cada artículo se extrajeron los siguientes elementos:
\begin{itemize}
    \item Objetivo del estudio
    \item Metodología empleada
    \item Conclusiones principales
    \item Limitaciones
\end{itemize}

A partir de esta información se construyó una matriz de antecedentes que resume los aportes más relevantes en el estudio de accidentes nucleares.

{\footnotesize
\renewcommand{\arraystretch}{1.2}
\begin{longtable}{|p{2.2cm}|p{0.7cm}|p{2.4cm}|p{2.8cm}|p{2.8cm}|p{2.8cm}|p{2.3cm}|}
\hline
\textbf{Autor} & \textbf{Año} & \textbf{Título} & \textbf{Objetivo} & \textbf{Metodología} & \textbf{Resultado} & \textbf{Limitaciones} \\
\hline
\endfirsthead

\hline
\textbf{Autor} & \textbf{Año} & \textbf{Título} & \textbf{Objetivo} & \textbf{Metodología} & \textbf{Resultado} & \textbf{Limitaciones} \\
\hline
\endhead

Y.H. Koo; Y.S. Yang; K.W. Song 
& 2014 
& Radioactivity release from the Fukushima accident and its consequences: A review.
& Analizar el impacto del accidente de Fukushima en la liberación de radiactividad y sus consecuencias a nivel ambiental, sanitario y energético.
& Análisis estadístico de series temporales combinado con modelos de dispersión atmosférica.
& Se evidencia un impacto global significativo de la liberación radiactiva, con alta incertidumbre en los efectos a largo plazo de la exposición a bajas dosis.
& Alta variabilidad en los resultados y falta de consenso en la estimación de efectos a largo plazo. \\
\hline


J. Lelieveld; D. Kunkel; M.G. Lawrence 
& 2011 
& Global risk of radioactive fallout afterM major nuclear reactor accidents.
& Evaluar el riesgo asociado a la lluvia radiactiva tras accidentes nucleares severos (nivel INES 7).
& Modelos de dispersión atmosférica basados en datos del accidente de Chernóbil, junto con modelos de desintegración nuclear.
& Se evidencia que los riesgos asociados a accidentes nucleares pueden ser significativamente mayores a los estimados en décadas anteriores.
& Alta incertidumbre asociada a la limitada disponibilidad de datos empíricos, la simplificación de variables ambientales y el enfoque en accidentes de gran escala. \\
\hline

E. de la Cruz-Sánchez; J. Klapp; E. Mayoral-Villa; R. González-Galán; A.M. Gómez-Torres
& 2018
& Numerical simulation of a temporary repository of radioactive material.
& Utilizar técnicas de simulación computacional para evaluar y seleccionar el sitio más apropiado para el aislamiento de radionuclidos.
& Modelamiento utilizando el Método de Elementos Finitos. Se emplea la ley de Darcy para calcular el campo de locidad, para resolver la ecuación de transporte de masa.
& Se desarrolló un modelo validado capaz de describir con precisión el movimiento de radionúclidos en medios porosos saturados con presencia de fallas.
& El modelo es bidimensional y podría no capturar la complejidad total de un sistema tridimensional real. \\
\hline

V. Kashparov; B. Salbu; S. Levchuk; V. Protsak; I. Maloshtan; C. Simonucci; C. Courbet; H.L. Nguyen; N. Sanzharova; V. Zabrotsky
& 2019
& Environmental behaviour of radioactive particles from Chernobyl
& Predecir la tasa de liberación del \ce{^{90}Sr} en campos de cultivo, trincheras de desechos y en el estanque de la planta de Chernóbil.
& Extracciones secuenciales y técnicas de radiografía.
& El pico de movilidad del \ce{^{90}Sr} ya ocurrió, lo que sugiere una disminución progresiva del riesgo en la cadena alimentaria. Además, una fracción significativa del material en las trincheras se encuentra en formas químicas altamente estables.
& Contaminación no uniforme y dificultad en la obtención y procesamiento de datos representativos para todo el terreno. \\
\hline

S. Sugihara
& 2013
& Deactivation of radiation from radioactive materials contaminated in a nuclear power plant accident
& Validar la teoría de que la interacción entre particulas generadas por agua procesada (infotones) y los núcleos de cesio puede transformarlos en elementos estables.
& Se toma muestras de suelos radioactivos para comprobar el cambio a elementos estables en presencia de agua activada.
& El tratamiento con agua activada que contiene infotones logró una reducción drástica de la radioactividad.
& Los infotones son partículas que aun no se pueden detectar directamente y no hay forma cualitativa de explicar el proceso estudiado. \\
\hline

Isotope Shifts, Radioactive Decay, and the Nuclear Forces
& W. Raven
& 2024
& (Book chapter / not specified)
& This chapter explores nuclear structure, focusing on how the number of neutrons affects atomic transition frequencies and nuclear stability. It discusses the role of nuclear forces in governing stability and radioactive decay. Various decay modes are analyzed, including neutron emission, proton emission, $\alpha$ decay, spontaneous fission, $\beta^-$ decay, $\beta^+$ decay, and electron capture.\\
\hline

Radioactive Decay 
& Paola Cappellaro 
& 2023 
& Physics LibreTexts (MIT OpenCourseWare, educational resource) 
& Describes radioactive decay as the process by which an unstable nucleus spontaneously loses energy by emitting particles and ionizing radiation; introduces the concept of nuclear instability and its spontaneous nature within the framework of basic nuclear physics. \\
\hline

\end{longtable}
}
